problem:  
Let \(M^n\) be a closed, simply connected Riemannian manifold with **nonnegative sectional curvature** and an **isometric, fixed point homogeneous action** by a compact connected Lie group \(G\). Let \(F\subset \mathrm{Fix}(G,M)\) be the maximal-dimensional fixed point component, and let \(Q = M/G\) be the corresponding **Alexandrov space** with curvature \(\ge 0\) and boundary \(\partial Q = F\).  

For each boundary point \(p \in F\), let \(L_p = \nu_p^1 F / G_p\) denote the **link**—the Alexandrov quotient of the unit normal sphere \(\nu_p^1 F\) by the isotropy representation of \(G_p\). These \(L_p\) are Alexandrov spaces of curvature \(\ge 1\) that describe the local cross-sectional geometry of \(Q\) along \(\partial Q\).  

**Problem:**  
Assume that, for every \(p \in F\), the link \(L_p\) is **isometric** (not merely homeomorphic) to a fixed compact Alexandrov space \(\Sigma\) of curvature \(1\).  

> Does it follow that the orbit space \(Q = M/G\) must admit a global **warped product decomposition**
> \[
>  (Q,d_Q) \cong ([0,L],dt^2) \times_\psi (\Sigma,d_{\Sigma}),
> \]
> for some concave warping function \(\psi:[0,L]\to \mathbb{R}_{>0}\) and some \(L>0\)?  
>
> Equivalently, does **constancy of the boundary link geometry** imply a **global radial product structure** for \(Q\) under nonnegative curvature and fixed point homogeneity?

If not, can one construct or preclude an example of a smooth nonnegatively curved manifold \(M\) with such a \(G\)-action for which all boundary links \(L_p\) are isometric but \(Q\) fails to admit any global warped (or even topological) product structure?  

---

Why is it a "good" problem:  
This problem directly probes a **new geometric rigidity mechanism** at the interface of **transformation group theory** and **Alexandrov geometry**.  

- **Conceptual essence:** It isolates a foundational question about how **local isotropy geometry** along the fixed set boundary governs the **global Alexandrov geometry** of the quotient. While fixed point homogeneity ensures local join or cone structures near \(\partial Q\), nothing is currently known about whether global warping follows purely from the uniformity of these infinitesimal links.  

- **Theoretical bridge:** The problem connects multiple fields:  
  - **Alexandrov geometry** (boundary structure, links, curvature propagation, concave retractions);  
  - **Riemannian group actions** (isotropy representations, slice geometry, orbit-type stratification);  
  - **Global analysis** (the PDE and convexity properties governing warped product curvatures).  
  Progress would require uniting these strands, possibly developing new **boundary-link rigidity theorems** in synthetic curvature settings.  

- **Analytic and topological depth:** Proving such rigidity—or building counterexamples—demands understanding whether uniform infinitesimal geometry constrains large-scale curvature collapse or branch behavior in orbit spaces, a delicate mix of **metric, topological, and group-theoretic data**.  

- **Simplicity with depth:** The assumption and desired conclusion are expressible in simple geometric terms—“all boundaries look the same; is the space radially homogeneous?”—yet answering it touches one of the most subtle questions in comparison geometry: *when does local curvature data globalize?*  

Hence, it is a “good” problem: it is cleanly formulated, lies exactly at a research frontier where none of the standard classification or rigidity results apply, and addressing it would yield new structural insight into the global metric geometry of nonnegatively curved, symmetric manifolds.